\section{Methodology}
\label{sec:methodology}

Building on our observation that UQ method effectiveness may depend on task format, we design a controlled comparison methodology that enables fair evaluation while isolating format-specific effects.

\subsection{Overview}

Our methodology has three components: (1) benchmark selection for ground-truth labels without annotation, (2) controlled variables ensuring fair comparison, and (3) mechanism verification separating implementation correctness from discrimination effectiveness.

\textbf{Rationale:} Prior comparisons conflated method differences with evaluation differences (different benchmarks, splits, models). Our design holds evaluation constant, varying only the UQ method under test.

\subsection{Benchmark Selection}

We use TruthfulQA mc1 (multiple-choice, single correct answer) as our primary benchmark.

\textbf{Rationale:} MC format provides ground-truth labels without human annotation---the model's selected answer is either correct or incorrect. This enables AUROC computation for hallucination detection where ``hallucination'' = incorrect answer selection. We acknowledge this scopes our results to MC format; free-form generation evaluation is future work.

\textbf{Dataset Statistics:}
\begin{itemize}
    \item Total questions: 817 (full), 50 (PoC validation)
    \item Format: 4 answer choices (A/B/C/D) per question
    \item Labels: Binary (correct / hallucination)
\end{itemize}

\subsection{Model Selection}

We use Llama-3-8B-Instruct as our primary evaluation model.

\textbf{Rationale:} Representative of the 7--13B decoder-only instruction-tuned model class. Open-weight availability enables reproducibility. Instruction tuning provides appropriate MC format handling.

\subsection{UQ Methods Under Test}

We evaluate three UQ methods spanning token-level and semantic-level approaches:

\textbf{Max Probability (Token-Level):} Uncertainty score:
\begin{equation}
u = 1 - \max_c P(c) \quad \text{where } c \in \{A, B, C, D\}
\end{equation}

\textbf{Choice Entropy (Token-Level):} Uncertainty score:
\begin{equation}
u = H(P) = -\sum_c P(c) \log P(c)
\end{equation}

\textbf{Semantic Entropy (Semantic-Level):} Uncertainty score:
\begin{equation}
u = H(P_{\text{cluster}})
\end{equation}
computed over NLI-clustered responses from $N=5$ samples.

\subsection{Mechanism Verification}

A key methodological contribution is separating mechanism verification from discrimination evaluation.

\textbf{Problem:} If semantic entropy achieves low AUROC, is the implementation broken or is the method unsuited to the task?

\textbf{Solution:} We verify that semantic clustering functions correctly independent of discrimination:
\begin{itemize}
    \item \textbf{Cluster count check:} Average clusters $<$ samples indicates clustering occurs
    \item \textbf{Entropy variance check:} Standard deviation of entropy $> 0$ indicates signal variation
\end{itemize}

If mechanism verification passes but discrimination fails, the method is correctly implemented but unsuited to the task format.
